\chapter{자주 쓰는 RTL 구조의 Area--Timing Trade-off}

\begin{summarybox}[좋은 코드는 목표에 따라 달라진다]
면적만 줄이는 코드와 timing만 줄이는 코드는 다를 수 있다. 이 장은 자주 쓰는
RTL 패턴을 생성 회로, area, timing, trade-off와 verification 순으로 비교한다.
Unified NTT 사례는 각 판단이 실제 설계에서 어떻게 나타났는지 확인하는 근거다.
\end{summarybox}

\section{긴 adder chain과 balanced tree}

여러 operand를 왼쪽부터 더하면 RTL은 짧아도 carry가 직렬로 이어질 수 있다.

\begin{lstlisting}
// A long dependency chain may be formed.
wire [24:0] sum_chain = a + b + c + d;
\end{lstlisting}

독립 operand라면 두 쌍을 먼저 더해 tree depth를 줄일 수 있다.

\begin{lstlisting}
wire [23:0] s0 = {1'b0, a} + {1'b0, b};
wire [23:0] s1 = {1'b0, c} + {1'b0, d};
wire [24:0] sum_tree = {1'b0, s0} + {1'b0, s1};
\end{lstlisting}

두 구조의 adder 수는 비슷하지만 dependency depth가 다르다. 합성기가 associative
expression을 재균형화할 수 있어도 signedness, truncation과 saturation이 끼면
자유롭게 바꾸기 어렵다.

\begin{itemize}
  \item \textbf{Area}: 총 carry/LUT 수는 비슷할 수 있고 intermediate guard bit가 늘 수 있다.
  \item \textbf{Timing}: chain은 carry depth가 세 번, balanced tree는 두 level에서 시작한다.
  \item \textbf{Verification}: 각 intermediate 폭과 최종 truncation이 원식과 같은지 확인한다.
\end{itemize}

Montgomery 경로에서 여러 9-bit operand를 한 식에 모았던 실험은 carry와 MUX가
깊게 결합됐다. 최종판은 일부 합을 앞 stage에서 만들고 다음 stage에서 남은 항을
더해 cycle 경계를 사용했다.

\section{MUX-then-arithmetic와 arithmetic-then-MUX}

\begin{lstlisting}
// Operator sharing favors area.
wire [22:0] y_shared = a + (mode ? b : c);

// Parallel arithmetic can shorten the select-to-result path.
wire [22:0] y_parallel = mode ? (a + b) : (a + c);
\end{lstlisting}

\begin{itemize}
  \item \textbf{Area}: 첫 식은 operand MUX와 adder 하나, 두 번째는 adder 두 개와 result MUX 후보이다.
  \item \textbf{Timing}: 첫 식은 MUX 뒤에 carry가 시작되고, 두 번째는 add와 select가 병렬 진행될 수 있다.
  \item \textbf{Trade-off}: DSP처럼 큰 operator는 공유 가치가 크고 작은 adder는 MUX 비용이 상대적으로 크다.
  \item \textbf{Verification}: mode가 data와 같은 stage인지 확인한다.
\end{itemize}

최종 \rtlfile{unifiedMUL}은 DSP 두 개를 여러 mode가 공유하는 쪽을 선택했다. 반면
butterfly 첫 stage의 mode별 결과를 너무 빨리 공통 rail로 합치면 carry 뒤 result
MUX가 critical path에 붙었으므로 결과를 한 stage 더 분리했다.

\section{Compare--subtract--correction chain}

\begin{lstlisting}
wire [22:0] r0 = (x >= Q) ? x - Q : x;
wire [22:0] r1 = (r0 >= Q) ? r0 - Q : r0;
\end{lstlisting}

짧은 두 줄이지만 comparator, subtractor와 MUX가 반복된다. 입력 범위를 증명하면
보정 횟수를 줄이거나 raw subtract의 sign/borrow를 직접 사용할 수 있다.

\begin{lstlisting}
wire [23:0] raw = {1'b0, a} - {1'b0, b};
wire [23:0] addq = raw + {1'b0, Q};
wire [22:0] y = raw[23] ? addq[22:0] : raw[22:0];
\end{lstlisting}

이 표현도 carry 두 번을 한 cycle에 포함할 수 있다. 150 MHz가 목표라면 raw와
correction 사이 register를 검토한다. range proof가 없는 correction 제거는
면적 최적화가 아니라 기능 오류다.

\section{Variable shift와 enumerated constant shift}

\begin{lstlisting}
wire [22:0] y_var = x << shamt;
\end{lstlisting}

runtime \rtlfile{shamt}는 여러 shifted value를 고르는 barrel MUX를 만든다. 실제
지원 값이 세 개뿐이라면 세 constant wiring 후보를 case로 선택할 수 있다.

\begin{lstlisting}
always @* begin
  case (mode)
    2'd0: y = x << 4;
    2'd1: y = x << 6;
    default: y = x << 7;
  endcase
end
\end{lstlisting}

이 경우 shift는 wiring이고 mode MUX만 남는다. 이번 23-bit 실험에서 constant
shift는 LUT 0, variable shift는 LUT 64였다. 지원하지 않는 shift 값을 제거한
것이므로 mode 범위 assertion과 exhaustive test가 필요하다.

\section{Priority chain과 parallel selection}

\begin{lstlisting}
if      (req[0]) y = d0;
else if (req[1]) y = d1;
else if (req[2]) y = d2;
else             y = d3;
\end{lstlisting}

여러 request가 동시에 1일 때 낮은 index가 이겨야 한다면 priority chain이 기능의
일부다. 원래 one-hot selector라면 case/one-hot MUX와 assertion으로 의도를 분리할
수 있다.

\begin{lstlisting}
case (sel)
  2'd0: y = d0;
  2'd1: y = d1;
  2'd2: y = d2;
  default: y = d3;
endcase
\end{lstlisting}

\begin{itemize}
  \item priority chain은 조건 수가 늘수록 select dependency가 길어질 수 있다.
  \item parallel selection은 decode fanout과 data MUX가 생긴다.
  \item 기능이 다른 두 구조를 timing만 보고 바꾸지 않는다.
\end{itemize}

\section{Decode-first와 arithmetic control}

복잡한 mode 비교를 여러 arithmetic branch마다 반복하면 decode logic와 high-fanout
net이 생길 수 있다.

\begin{lstlisting}
wire is_r3 = (mode == MODE_R3_A) || (mode == MODE_R3_B);
reg  is_r3_q;
always @(posedge clk)
  is_r3_q <= is_r3;

wire [22:0] y = is_r3_q ? r3_value : r2_value;
\end{lstlisting}

decode를 한 번 계산해 register하면 다음 stage는 1-bit control만 본다. FF 한
bit와 latency가 늘지만 반복 comparator와 긴 control route를 줄일 수 있다.
같은 mode가 transaction 동안 고정된다는 전제가 있다면 stage마다 decode를 다시
계산하지 않는 구조도 가능하다.

\section{Register replication과 high fanout}

Fanout은 하나의 신호가 구동하는 input pin의 수다. 아래처럼 하나의 enable을
datapath register, BRAM control과 valid 생성에 함께 쓰면 lane과 port를 추가할
때마다 load가 늘어난다.

\begin{lstlisting}
wire core_enable = (state == ST_CORE) && run;

always @(posedge clk) begin
  if (core_enable) begin
    a_pipe <= a_in;
    b_pipe <= b_in;
    c_pipe <= c_in;
    d_pipe <= d_in;
  end
end

assign bram_a_en = core_enable && bank_sel_a;
assign bram_b_en = core_enable && bank_sel_b;
assign valid_in  = core_enable && input_ready;
\end{lstlisting}

RTL만 보면 \rtlfile{core\_enable}은 비교와 AND 몇 개로 만든 작은 신호다. 하지만
합성 후에는 여러 LUT의 select와 여러 FF의 enable pin까지 연결된다. sink가 서로
멀리 배치되면 하나의 net이 넓은 영역을 지나가므로 logic가 작아도 route delay가
critical path를 지배할 수 있다.

내 200~MHz 목표 RTL의 LOAD-to-BRAM 경로에서 이 문제가 실제로 나타났다.

\begin{center}
\begin{tabular}{lr}
\toprule
항목 & 측정값 \\
\midrule
전체 data-path delay & 7.913~ns \\
Logic delay & 1.448~ns \\
Route delay & 6.465~ns \\
Route 비율 & 81.7\% \\
Logic level & LUT 8단 \\
관찰된 fanout & 95 \\
\bottomrule
\end{tabular}
\end{center}

처음에는 LUT 8단을 먼저 줄이고 싶었지만, 전체 지연의 81.7\%가 배선이었다.
이 경우 Boolean 식을 조금 줄이는 것만으로 5 ns 안에 들어가기 어렵다. LOAD 상태,
layout mapping과 memory control이 너무 많은 BRAM input까지 퍼진 구조가 핵심이었다.

다음처럼 wire 이름만 나누어도 실제 driver는 복제되지 않는다.

\begin{lstlisting}
wire load_en_a = load_enable;
wire load_en_b = load_enable;
wire load_en_c = load_enable;
\end{lstlisting}

합성기는 세 alias가 같은 값임을 알고 하나의 net으로 합칠 수 있다. 내가 먼저 한
일은 register를 무조건 복제하는 것이 아니라, LOAD 요청을 한 번 physical memory
command로 변환하고 BRAM-facing boundary에 저장하는 것이었다.

\begin{lstlisting}
always @(posedge clk) begin
  if (load_accept) begin
    host_valid_q <= 1'b1;
    host_bank_q  <= mapped_bank;
    host_addr_q  <= mapped_addr;
    host_data_q  <= load_data;
    host_we_q    <= mapped_we;
  end else begin
    host_valid_q <= 1'b0;
  end
end

assign bram_addr = host_addr_q;
assign bram_we   = host_valid_q && host_we_q;
assign bram_din  = host_data_q;
\end{lstlisting}

이 구조에서는 LOAD/BUSY/layout 조건이 BRAM port마다 다시 계산되지 않는다.
address, bank와 write enable도 같은 register 경계에서 정렬된다. 최종 150~MHz
구현에서는 LOAD/DUMP address path가 worst path에서 사라지고 병목이 butterfly
arithmetic으로 이동했다. fanout 하나의 최종 숫자만 비교한 결과는 아니지만,
high-fanout control을 BRAM-facing cone에서 제거한 효과를 path 이동으로 확인했다.

그래도 하나의 enable register가 물리적으로 먼 여러 영역을 계속 구동하면
region별 register replication을 검토한다.

\begin{lstlisting}
(* keep = "true" *)
reg en_left_q, en_right_q;
always @(posedge clk) begin
  en_left_q  <= en;
  en_right_q <= en;
end
\end{lstlisting}

왼쪽 datapath는 \rtlfile{en\_left\_q}, 오른쪽 datapath는
\rtlfile{en\_right\_q}를 사용한다. 기능적으로 같은 register를 region별로
복제하면 load와 물리적 영역을 나눌 수 있다.

\begin{itemize}
  \item \textbf{Area}: FF와 clock load가 조금 증가한다.
  \item \textbf{Timing}: control route와 skewed arrival가 줄 수 있다.
  \item \textbf{Tool interaction}: 합성기가 equivalent register를 merge할 수 있다.
  \item \textbf{Verification}: 두 복제 값이 항상 같고 reset/enable도 일치해야 한다.
\end{itemize}

복제된 register는 합성기가 다시 합칠 수 있으므로 실제 netlist에서 cell이
유지됐는지 확인한다. 내가 fanout 문제를 보는 순서는 다음과 같다.

\begin{enumerate}
  \item fanout 숫자와 함께 route-delay 비율을 확인한다.
  \item 그 control이 필요하지 않은 block까지 퍼지는지 찾는다.
  \item 반복되는 state/mode 해석을 local registered command로 바꾼다.
  \item 그래도 물리적으로 넓으면 region별 register replication을 검토한다.
  \item 수정 후 새 worst path가 어디로 이동했는지 확인한다.
\end{enumerate}

\section{Pipeline placement: 식 가운데가 아니라 경로 경계}

\begin{lstlisting}
// One long stage
always @(posedge clk)
  y_q <= ((a + b) >= Q) ? (a + b - Q) : (a + b);
\end{lstlisting}

raw sum을 register하면 carry와 correction을 다른 cycle에 둘 수 있다.

\begin{lstlisting}
always @(posedge clk) begin
  raw_q <= {1'b0, a} + {1'b0, b};
  y_q   <= (raw_q >= Q) ? raw_q - Q : raw_q;
end
\end{lstlisting}

\begin{itemize}
  \item \textbf{Area}: raw data와 valid/mode/address를 위한 FF가 늘어난다.
  \item \textbf{Timing}: carry chain 사이에 register boundary가 생긴다.
  \item \textbf{Latency}: 한 transaction의 결과가 한 cycle 늦어진다.
  \item \textbf{Throughput}: 두 stage가 매 cycle 동작하면 II=1을 유지할 수 있다.
\end{itemize}

NTRU+ constant-multiply path에 stage 하나를 추가했던 첫 시도는 arithmetic timing은
나아졌지만 다른 butterfly branch와 writeback address가 어긋나 기능이 깨졌다.
최종 수정은 모든 branch의 common delay와 schedule compensation을 함께 바꿨다.

\section{Minimum width와 guard bit}

모든 intermediate를 편의상 32 bit로 만들면 MUX, register, adder와 comparator가
그 폭으로 반복될 수 있다. 실제 범위가 9 bit라면 9-bit datapath로 유지하는 것이
area와 carry depth에 유리하다.

반대로 modular add/subtract의 carry, borrow 또는 signed correction을 위해 한 bit
guard가 필요할 수 있다. 최소 폭은 ``가장 작은 선언''이 아니라 모든 가능한
intermediate를 보존하는 가장 작은 폭이다.

\begin{enumerate}
  \item operand range를 적는다.
  \item add/subtract/multiply 뒤의 최대·최소값을 계산한다.
  \item 필요한 signed bit와 guard bit를 정한다.
  \item truncation 시점과 modular equivalence를 증명한다.
  \item boundary value를 simulation/formal assertion으로 확인한다.
\end{enumerate}

\section{Reset과 enable을 필요한 state에만 둔다}

data pipeline 전체를 reset하면 control set과 routing이 늘고 SRL/DSP/BRAM register
inference를 막을 수 있다. valid가 0일 때 data가 관찰되지 않는 구조라면 valid와
protocol state만 reset할 수 있다.

이 최적화는 initialization을 무시하는 방법이 아니다. reset release 뒤 first
valid transaction, backpressure, pipeline flush와 stale data가 output에 노출되지
않는다는 계약을 검증해야 한다.

\section{패턴을 적용하기 전의 질문}

\begin{checklistbox}
\begin{enumerate}
  \item 이 변경은 LUT, FF, DSP와 BRAM 중 무엇을 줄이고 무엇을 늘리는가?
  \item critical path에서 MUX, carry와 route의 순서가 어떻게 바뀌는가?
  \item latency와 initiation interval이 각각 어떻게 변하는가?
  \item data와 valid, mode, address의 alignment를 어디서 보장하는가?
  \item 합성기가 다시 공유, 복제 또는 trimming할 가능성을 report에서 확인했는가?
  \item 같은 top/device/constraint의 변경 전후 결과가 있는가?
\end{enumerate}
\end{checklistbox}
